\documentclass[11pt,a4paper]{article}
\usepackage[utf8]{inputenc}
\usepackage{amsmath,amssymb,amsthm}
\usepackage[margin=1in]{geometry}

\newtheorem{theorem}{Theorem}
\newtheorem{lemma}[theorem]{Lemma}
\newtheorem{definition}[theorem]{Definition}

\title{Problem 124: Ordered Radicals}
\author{Project Euler Solution}
\date{}

\begin{document}
\maketitle

\section{Problem Statement}

Define $\operatorname{rad}(n) = \prod_{p \mid n} p$ (product of distinct prime factors), with $\operatorname{rad}(1) = 1$. Sort the pairs $(\operatorname{rad}(n), n)$ for $1 \le n \le 100{,}000$ lexicographically. Find $E(10000)$.

\section{Mathematical Development}

\begin{definition}
For $n = p_1^{e_1} \cdots p_m^{e_m}$, the radical is $\operatorname{rad}(n) = p_1 \cdots p_m$.
\end{definition}

\begin{theorem}[Squarefreeness characterization]
$\operatorname{rad}(n) = n$ if and only if $n$ is squarefree. For any prime power, $\operatorname{rad}(p^k) = p$.
\end{theorem}

\begin{proof}
If $n = p_1 \cdots p_m$ is squarefree, then $\operatorname{rad}(n) = n$. If $p^2 \mid n$, then $n > \operatorname{rad}(n)$. For $p^k$, the unique prime factor is $p$.
\end{proof}

\begin{theorem}[Sieve correctness]
Initialize $\operatorname{rad}[n] = 1$. For each prime $p$, multiply $\operatorname{rad}[m]$ by $p$ for all multiples $m$ of $p$. This yields $\operatorname{rad}[n] = \prod_{p \mid n} p$.
\end{theorem}

\begin{proof}
Each prime $p \mid n$ visits $n$ exactly once, contributing one factor of $p$. Primes act independently and are identified as those $p$ with $\operatorname{rad}[p] = 1$ upon first encounter.
\end{proof}

\begin{lemma}[Total order]
The lexicographic order on $(\operatorname{rad}(n), n)$ is total: $n$ is a unique tiebreaker.
\end{lemma}

\begin{proof}
Immediate from the uniqueness of $n$.
\end{proof}

\section{Algorithm}

\begin{enumerate}
\item Sieve $\operatorname{rad}[n]$ for $n = 1, \ldots, N$.
\item Sort the integers $1, \ldots, N$ by $(\operatorname{rad}(n), n)$.
\item Return the $n$-value at position $K = 10{,}000$.
\end{enumerate}

\section{Complexity Analysis}

\begin{itemize}
\item \textbf{Time:} $O(N \log \log N)$ for sieve (by Mertens' theorem), $O(N \log N)$ for sort.
\item \textbf{Space:} $O(N)$.
\end{itemize}

\section{Answer}

\[
\boxed{21417}
\]

\end{document}
